\documentclass[authoryear,preprint,review,12pt]{elsarticle}

\usepackage[utf8]{inputenc}
\usepackage[T1]{fontenc}
\usepackage{amsmath}
\usepackage{amssymb}
\usepackage{graphicx}
\usepackage{hyperref}
\usepackage{booktabs}
\usepackage{xcolor}
\usepackage{listings}
\usepackage{url}
\usepackage{float}
\usepackage{subcaption}

% Code listing style
\lstset{
  basicstyle=\ttfamily\small,
  breaklines=true,
  frame=single,
  language=R,
  keywordstyle=\color{blue},
  commentstyle=\color{gray},
  stringstyle=\color{red},
  numbers=left,
  numberstyle=\tiny\color{gray},
  numbersep=5pt
}

% Path to figures
\graphicspath{{../output/figures/}}

\journal{SoftwareX}

\begin{document}

\begin{frontmatter}

\title{crossover-edtech-toolkit: An Open-Source Platform for Replicable\\
Crossover Experimental Studies Evaluating Generative AI in Education}

\author[uex]{Jes\'us \'Angel Torrecilla-Pinero\corref{cor}}
\ead{jatorrecilla@unex.es}
\cortext[cor]{Corresponding author}

% Add remaining authors here
%\author[uex]{Second Author}
%\author[uex]{Third Author}

\address[uex]{Department of Construction, Escuela Polit\'ecnica, Universidad de Extremadura, C\'aceres, Spain}

\begin{abstract}
Evaluating the impact of generative artificial intelligence (AI) on learning outcomes
requires rigorous experimental designs that control for individual differences among
students. This paper presents \texttt{crossover-edtech-toolkit}, an open-source
software platform that provides a complete framework for conducting AB/BA crossover
experimental studies in higher education. The toolkit comprises four integrated
components: (1)~a web-based data collection application built on Firebase with
role-based access, anonymous data coding, and conditional survey logic;
(2)~a set of validated measurement instruments including rubrics and Likert-scale
questionnaires with documented psychometric properties (Aiken's $V = 0.65$--$0.95$,
Cronbach's $\alpha > 0.70$);
(3)~dual analysis pipelines in R (10 modular scripts) and Python (10 scripts plus
Jupyter notebook) implementing the complete statistical workflow for 2$\times$2
crossover designs, from data cleaning through mixed ANOVA, linear mixed models,
Grizzle's carryover test, and publication-quality visualizations; and
(4)~comprehensive documentation with tutorials, study design guides, and adaptation
manuals for different disciplines. We demonstrate the toolkit's capabilities using
a synthetic dataset of 100 students, reporting a significant treatment effect
(Cohen's $d = 0.43$, $p < 0.001$), non-significant carryover ($p = 0.14$),
and consistent linear mixed model estimates. The software is released under the
MIT license and is available on GitHub with a citable DOI.
\end{abstract}

\begin{keyword}
crossover design \sep educational technology \sep generative AI \sep
higher education \sep open-source software \sep experimental methodology \sep
replicability \sep statistical analysis pipeline
\end{keyword}

\end{frontmatter}

% ==============================================================================
\section{Motivation and significance}
\label{sec:motivation}
% ==============================================================================

The rapid adoption of generative AI tools (e.g., ChatGPT, GitHub Copilot, Gemini,
Claude) in higher education has created an urgent need for rigorous empirical
evidence about their effects on learning outcomes \citep{kasneci2023chatgpt}.
While numerous studies have explored student and instructor perceptions of AI
tools, far fewer have employed experimental designs capable of establishing
causal claims \citep{crompton2023ai, yan2024practical}.

A key challenge in educational experiments is the large variability between
students in terms of prior knowledge, motivation, and ability. Traditional
parallel-group designs (treatment vs.\ control) require large sample sizes to
achieve adequate statistical power. The \textbf{crossover design}, widely used
in clinical trials \citep{senn2002crossover}, addresses this by having each
participant experience both conditions: with AI and without AI. This
within-subject comparison eliminates between-subject variability, yielding
greater statistical power with smaller samples \citep{wellek2012crossover,
jones2014crossover}.

Despite these advantages, crossover designs remain underutilized in educational
research, partly because there is no standardized, freely available toolkit for
implementing them. Researchers must build data collection platforms, design
instruments, and write analysis code from scratch, leading to methodological
inconsistencies and limited replicability \citep{nosek2015promoting}.

\texttt{crossover-edtech-toolkit} addresses this gap by providing a complete,
open-source platform that can be adopted by any institution or discipline.
The toolkit lowers the barrier to conducting high-quality crossover experiments
and promotes open science practices in educational technology research.

% ==============================================================================
\section{Software description}
\label{sec:description}
% ==============================================================================

\subsection{Software architecture}

The toolkit follows a modular architecture with four loosely coupled components
(Figure~\ref{fig:architecture}):

\begin{enumerate}
    \item \textbf{Data collection webapp}: A Firebase-hosted web application
    that guides participants through the study protocol.

    \item \textbf{Validated instruments}: A set of questionnaires, rubrics,
    and surveys with documented psychometric properties.

    \item \textbf{R analysis pipeline}: Ten modular R scripts implementing
    the complete statistical analysis for a 2$\times$2 crossover design.

    \item \textbf{Python analysis pipeline}: Ten equivalent Python scripts
    plus an interactive Jupyter notebook, producing identical statistical
    results using pandas, scipy, pingouin, and statsmodels.
\end{enumerate}

\begin{figure}[htbp]
\centering
\includegraphics[width=\textwidth]{architecture.pdf}
\caption{Software architecture of the \texttt{crossover-edtech-toolkit}. Three data-producing components (web application, validated instruments, and sample data generator) feed into dual analysis pipelines (R and Python) that produce a standardized set of outputs (tables, figures, fitted models, and execution logs) ready for publication. Comprehensive documentation supports all stages.}
\label{fig:architecture}
\end{figure}

\subsection{Software functionalities}

\subsubsection{Data collection}

The Firebase webapp (\url{https://gestor-investigacion-app.web.app})
implements the following workflow:

\begin{enumerate}
    \item Participant registration via magic link (passwordless authentication)
    \item Random assignment to sequence AB or BA
    \item Pre-test knowledge and attitude assessment
    \item Period 1 challenge task (condition determined by sequence)
    \item Post-challenge perception survey (conditional items based on condition)
    \item Period 2 challenge task (alternate condition)
    \item Post-challenge perception survey
    \item Post-test reflective questionnaire
\end{enumerate}

Key technical features of the webapp include:
\begin{itemize}
    \item \textbf{Role-based access}: Three user roles (student, teacher, admin)
    with appropriate permissions enforced via Firestore security rules.
    \item \textbf{Anonymous data coding}: Student identifiers are hashed using
    SHA-256, preventing identification by researchers while maintaining data
    linkability across time points.
    \item \textbf{Conditional survey logic}: Post-challenge items adapt to the
    experimental condition (AI vs.\ no-AI), collecting AI-specific data
    (prompt strategies, tool usage patterns) only when applicable.
    \item \textbf{Teacher-blind evaluation}: Evaluators access submissions
    without knowledge of the student's condition or sequence.
    \item \textbf{Data export}: CSV/JSON export with automatic wide-format
    restructuring for analysis pipeline ingestion.
\end{itemize}

\subsubsection{Analysis pipelines}

Both the R and Python pipelines implement the standard analysis strategy for
crossover designs \citep{jones2014crossover}, organized as ten sequential
modules:

\begin{itemize}
    \item \textbf{00\_setup}: Package management and configuration
    \item \textbf{01\_data\_import\_clean}: Data validation, missing data
    handling, derived variables (within-subject differences, period sums)
    \item \textbf{02\_descriptive\_stats}: Summary statistics by condition,
    period, and sequence
    \item \textbf{03\_instrument\_validation}: Cronbach's $\alpha$ (with
    item-total statistics), Aiken's $V$, ICC
    \item \textbf{04\_paired\_comparisons}: Paired $t$-tests and Wilcoxon
    signed-rank tests for within-subject treatment comparisons
    \item \textbf{05\_mixed\_anova}: Mixed ANOVA and linear mixed models (LMM)
    \item \textbf{06\_effect\_sizes}: Cohen's $d$, Hedges' $g$, partial
    $\eta^2$ with confidence intervals
    \item \textbf{07\_carryover\_analysis}: Grizzle's test \citep{grizzle1965two}
    for carryover effects (parametric and non-parametric)
    \item \textbf{08\_perception\_analysis}: Likert scale analysis with
    Wilcoxon tests and diverging stacked bar charts
    \item \textbf{09\_visualizations}: Interaction plots, spaghetti plots,
    forest plots, composite multi-panel figures
\end{itemize}

The LMM follows the standard crossover parameterization:
\begin{equation}
    Y_{ijk} = \mu + \pi_j + \tau_k + \lambda_l + S_i + \varepsilon_{ijk}
    \label{eq:model}
\end{equation}
where $\mu$ is the grand mean, $\pi_j$ is the period effect, $\tau_k$ is
the treatment effect, $\lambda_l$ is the carryover effect, $S_i \sim
\mathcal{N}(0, \sigma_S^2)$ is the random subject effect, and
$\varepsilon_{ijk} \sim \mathcal{N}(0, \sigma^2)$ is the residual.

Both pipelines are orchestrated by a master script (\texttt{run\_all.R} or
\texttt{run\_all.py}) that executes all modules sequentially and logs
execution metadata (timestamps, package versions, row counts) to a pipeline
log file, supporting full reproducibility.

\subsubsection{Sample data generator}

The toolkit includes data generators in both R and Python that create
synthetic datasets with configurable parameters: sample size, treatment
effect, period effect, carryover effect, and variance components. The
default configuration simulates 100 students in a balanced 2$\times$2
crossover design. This generator enables researchers to:
\begin{itemize}
    \item Verify the pipeline installation before collecting real data
    \item Conduct power analyses via simulation
    \item Explore the behavior of the analytical pipeline under different
    scenarios (e.g., strong carryover, no treatment effect)
\end{itemize}

% ==============================================================================
\section{Illustrative examples}
\label{sec:examples}
% ==============================================================================

To demonstrate the toolkit's capabilities, we generated a synthetic dataset
of 100 students in a 2$\times$2 crossover design using the included data
generator. The data-generating parameters were: $\mu = 65$, $\tau = 5$
(AI advantage), $\pi = 3$ (period effect), $\lambda = 0$ (no carryover),
$\sigma_S = 12$, $\sigma = 8$.

Running the full R pipeline with \texttt{source("analysis/R/run\_all.R")}
produces 15 statistical tables, 15 publication-quality figures, and fitted
model objects, all organized in the \texttt{output/} directory.

\subsection{Key statistical results}

Table~\ref{tab:descriptive} shows descriptive statistics by condition.
Students scored an average of 5.19 points higher under the AI condition
($M = 71.34$ vs.\ $M = 66.15$) and required less time ($M = 35.82$ vs.\
$M = 41.88$ minutes).

\begin{table}[htbp]
\centering
\caption{Descriptive statistics by condition (synthetic data, $n = 100$).}
\label{tab:descriptive}
\begin{tabular}{lcccccc}
\toprule
\textbf{Condition} & $n$ & \textbf{Mean} & \textbf{SD} & \textbf{Median} & \textbf{Time (min)} & \textbf{SD time} \\
\midrule
No AI & 100 & 66.15 & 14.94 & 64.95 & 41.88 & 10.00 \\
AI    & 100 & 71.34 & 14.88 & 73.20 & 35.82 &  8.77 \\
\bottomrule
\end{tabular}
\end{table}

The carryover test was non-significant (Grizzle's $t = 1.48$, $p = 0.143$;
Wilcoxon $W = 1484$, $p = 0.107$), confirming the crossover comparison is
valid. The paired $t$-test showed a significant treatment effect ($t(99) =
4.27$, $p < 0.001$, 95\% CI $[2.78, 7.60]$), with Cohen's $d = 0.43$
[0.22, 0.63].

The mixed ANOVA confirmed the significant treatment effect ($F(1,98) = 18.85$,
$p < 0.001$, $\eta^2_G = 0.030$) and a significant sequence $\times$
condition interaction ($F(1,98) = 4.37$, $p = 0.039$), suggesting that
students who used AI second (Sequence AB) showed a larger effect ($d = 0.58$)
than those who used it first ($d = 0.26$). The LMM fixed effects were
consistent with the data-generating parameters: AI effect $= 5.19$ [2.85,
7.53], period effect $= 2.50$ [0.16, 4.84].

\subsection{Publication-quality visualizations}

The visualization module generates several figures designed for direct
inclusion in publications. Figure~\ref{fig:interaction} shows the interaction
between treatment and period by sequence---the crossing pattern is the
hallmark of a successful crossover detecting a treatment effect.

\begin{figure}[htbp]
\centering
\includegraphics[width=0.8\textwidth]{interaction_plot.png}
\caption{Interaction plot showing mean scores by sequence (AB vs.\ BA)
and period. The crossing pattern indicates that scores are higher in whichever
period includes AI assistance, regardless of sequence.}
\label{fig:interaction}
\end{figure}

Figure~\ref{fig:spaghetti} displays individual student trajectories across
periods, providing insight into the within-subject variability that the
crossover design leverages.

\begin{figure}[htbp]
\centering
\includegraphics[width=0.8\textwidth]{spaghetti_plot.png}
\caption{Spaghetti plot of individual student trajectories across periods,
colored by sequence. Each line represents one student.}
\label{fig:spaghetti}
\end{figure}

Figure~\ref{fig:forest} presents the effect sizes with confidence intervals,
enabling readers to assess both the magnitude and precision of the treatment
effect overall and by subgroup.

\begin{figure}[htbp]
\centering
\includegraphics[width=0.8\textwidth]{forest_plot.png}
\caption{Forest plot of effect sizes (Cohen's $d$) with 95\% confidence
intervals for the overall comparison and by sequence.}
\label{fig:forest}
\end{figure}

Figure~\ref{fig:composite} shows a composite multi-panel figure that
combines the key visualizations into a single publication-ready graphic.

\begin{figure}[htbp]
\centering
\includegraphics[width=0.95\textwidth]{composite_figure.png}
\caption{Composite figure combining (a) score distributions by condition,
(b) interaction by sequence and period, (c) individual trajectories, and
(d) effect size forest plot.}
\label{fig:composite}
\end{figure}

Additional figures generated by the pipeline include boxplots by condition
and period, violin plots, Q-Q plots for normality assessment, Bland-Altman
agreement plots, histograms of score differences, carryover diagnostic
plots, and Likert scale diverging bar charts (Figure~\ref{fig:likert}).

\begin{figure}[htbp]
\centering
\includegraphics[width=0.8\textwidth]{likert_comparison.png}
\caption{Comparison of Likert-scale perception responses between AI and
no-AI conditions.}
\label{fig:likert}
\end{figure}

% ==============================================================================
\section{Impact}
\label{sec:impact}
% ==============================================================================

The \texttt{crossover-edtech-toolkit} has the potential to impact educational
research in several ways:

\begin{enumerate}
    \item \textbf{Methodological improvement}: By providing a ready-to-use
    crossover framework, the toolkit encourages researchers to use
    within-subject designs that yield more precise estimates of treatment
    effects than parallel-group designs. A paired comparison with $n = 50$
    students in a crossover has comparable power to a parallel-group design
    with $n \approx 100$, making rigorous experiments feasible in typical
    university settings.

    \item \textbf{Replicability and comparability}: The standardized
    instruments, data format, and analysis pipeline enable direct
    replication across institutions and disciplines. When multiple sites
    use the same toolkit, results become directly comparable, enabling
    meta-analytic synthesis and multi-center studies.

    \item \textbf{Open science}: All components are open-source (MIT license),
    promoting transparency and enabling the research community to inspect,
    modify, and improve the methodology \citep{nosek2015promoting,
    wilson2017softwarex}.

    \item \textbf{Accessibility}: The comprehensive documentation---including
    a step-by-step tutorial, study design guide, and instrument adaptation
    manual---lowers the barrier for researchers without deep expertise in
    crossover designs or mixed-effects modeling. The dual R/Python
    implementation ensures accessibility for researchers from different
    statistical computing backgrounds.

    \item \textbf{Evidence base for AI policy}: As institutions develop
    policies on AI use in education, this toolkit provides a rigorous
    methodology for generating the empirical evidence needed to inform
    those decisions. The crossover design is particularly relevant because
    it produces within-subject evidence that is more convincing to
    policymakers than between-group comparisons.
\end{enumerate}

The toolkit is currently being used in a multi-course crossover study at the
Universidad de Extremadura (Spain), evaluating the impact of generative AI
on active learning in civil engineering programs. The study protocol,
including the use of this toolkit, is described in a companion methodological
paper.

% ==============================================================================
\section{Conclusions}
\label{sec:conclusions}
% ==============================================================================

We have presented \texttt{crossover-edtech-toolkit}, an open-source platform
for conducting crossover experimental studies evaluating generative AI in
higher education. The toolkit integrates data collection, validated
instruments, and dual-language statistical analysis into a coherent,
reusable framework. Using a synthetic dataset, we demonstrated the full
pipeline from data generation through publication-ready tables and
figures, confirming the implementation's correctness against known
data-generating parameters.

The toolkit addresses a critical gap in educational technology research:
the absence of standardized, accessible tools for implementing rigorous
within-subject experimental designs. By releasing all components under
the MIT license with comprehensive documentation, we aim to contribute
to a growing body of high-quality, replicable evidence on the educational
impact of AI tools.

Future development will include: (a) support for more complex crossover
designs (e.g., higher-order, replicated crossover with $> 2$ treatments);
(b) Bayesian analysis alternatives; (c) automated power analysis via
simulation; (d) multi-site federation features for coordinated data
collection across institutions; and (e) integration with learning
management systems (LMS) such as Moodle and Canvas.

% ==============================================================================
% Required metadata for SoftwareX
% ==============================================================================

\section*{Current code version}

\begin{table}[htbp]
\centering
\begin{tabular}{ll}
\toprule
\textbf{Code metadata} & \\
\midrule
Current code version & v1.0.0 \\
Permanent link to code/repository & \url{https://github.com/jtorreci/crossover-edtech-toolkit} \\
Permanent link to reproducible capsule & \url{https://doi.org/10.5281/zenodo.XXXXXXX} \\
Legal code license & MIT \\
Code versioning system used & Git \\
Software code languages & R, Python, JavaScript, HTML/CSS \\
Compilation requirements & R $\geq$ 4.2.0, Python $\geq$ 3.9, Node.js $\geq$ 18 \\
Developer documentation/manual & \url{https://jtorreci.github.io/crossover-edtech-toolkit/} \\
Support email for questions & jatorrecilla@unex.es \\
\bottomrule
\end{tabular}
\end{table}

% ==============================================================================
\section*{Acknowledgements}
% ==============================================================================

This work was supported by a Teaching Innovation Project (Proyecto de
Innovaci\'on Docente) at the Universidad de Extremadura (academic year
2025--2026). The authors thank the students who participated in the
pilot study and the colleagues in the teaching innovation group who
provided feedback on the instruments and analysis pipeline.

% ==============================================================================
\bibliographystyle{elsarticle-harv}
\bibliography{references}

\end{document}
